\chapter
[\CO{[SHF]} Sehen\,--\,Hören\,--\,Fühlen: Primary Survey I]
{\CC{[SHF]}\\Sehen\,--\,Hören\,--\,Fühlen:\\Basis-Patienteneinschätzung und Einführung in die "`Primary Survey"'}
\label{chp:BET}

\ChapterIntro
{%
---
}
{%
\begin{description}
  \item[Maintainer:] Sebastian Gabriel
  \item[Co-Maintainer:] Josef Emhofer, Roman Koch
  \item[Autoren:] Diverse
  \item[Reviewer:] Standard-Reviewprozess
  \item[Version:] {\DocVersion} ({\DocVersionInfo})
  \item[SHA1:] {\DocGitCommit}
\end{description}

{\TxTeilkapitelAASS}
}

\section{Sehen -- Hören -- Fühlen }
\label{topic:sehen-hoeren-fuehlen}
\Layout{0}{Unsere eigenen Sinne}{
Die Seh-, Hör- und Tastsinne sind bei der körperlichen
Untersuchung unverzichtbar. Zuerst wird der Patient und sein
Körper betrachtet. Wir versuchen, uns bereits bekannte Symptome wahrzunehmen.
Hierbei können uns viele augenscheinliche Sachen auffallen: zum Beispiel
\E{Fettleibigkeit}, \E{Zyanose}, \E{Blässe},
\E{Schweiß}, \E{Fremdkörper}, \E{Verletzungen}, \E{Fehlstellungen},
\ldots

Der nächste Schritt betrifft das Hören: Hier sind zum
Beispiel die Atemgeräusche einzuordnen, die man eventuell schon ohne
Hilfsmittel vernehmen kann (brodelndes Atemgeräusch beim
Lungenödem, pfeifendes Atemgeräusch beim Asthmaanfall,
röcheln bei Atemwegsverlegung, \ldots)

Wenn wir gesehen und gehört haben, können wir versuchen zu
fühlen: Wir können den (Radialis-) Puls fühlen, können
(stehende) Hautfalten beurteilen oder das Abdomen abtasten. Ebenso können wir
die ungefähre Temperatur fühlen. }{
\begin{itemize}
  \item Zyanose
  \item Blässe
  \item Schweiß
  \item Körperhaltung
  \item Fremdkörper
  \item Verletzungen, Fehlstellungen
  \item Atemgeräusche
  \item Puls
  \item Stehende Hautfalten
  \item Temperatur
  \item Durchblutung des Nagelbetts
  \item \ldots
\end{itemize}
}{}{}{}